
\def\PUDcodigo{ECA216}

\def\PUDcodacad{04507.29}

\def\PUDdisciplina{Eletrônica III}

\def\PUDsemestre{S6}

\def\PUDhoras{80}

%NOVOS ---------------------------------------------------
\def\PUDhorasTeoria{60}
\def\PUDhorasPratica{20}

\def\PUDinterECA{ 
\hyperref[PUD:ECA214]{Eletrônica II (04507.23)}

\hyperref[PUD:ECA215]{Circuitos Elétricos II (04507.24)}

\hyperref[PUD:ECA224]{Acionamento de Máquinas (04507.41)} 
}

\def\PUDatuaisECA{- Aplicações de eletrônica na Indústria 4.0

- Desenvolvimento de sistemas e circuitos eletrônicos de processamento de energia modernos}

\def\PUDrecursos{Projetor multimídia; Quadro branco e pincel; Aulas em laboratório de eletrônica.}

%----------------------------------------------------------

\def\PUDcreditos{4}

\def\PUDprerequisito{ \hyperref[PUD:ECA214]{ECA214} }

\def\PUDpreacad{ \hyperref[PUD:ECA214]{Eletrônica II (04507.23)} }

\def\PUDobjetivos{Compreender os principais elementos utilizados para controlar a potência elétrica, capacitando-se para o uso da eletrônica de potência na nas aplicações industriais e na Indústria 4.0. Compreender o processo de conversão de energia elétrica realizado pelos conversores CC-CC, aplicar os circuitos eletrônicos envolvidos no acionamento de interruptores de potência, compreender os processos de conversão de freqüência e de tensão nos conversores CA-CA. Familiarizar-se com a tecnologia e o desenvolvimento de sistemas e circuitos eletrônicos de processamento de energia modernos.}

\def\PUDementa{Introdução;  Interruptores Eletrônicos;  Cálculos de potência;  Retificadores de meia-onda;  Retificador de onda completa;  Retificadores trifásicos;  
Controladores de tensão CA (Gradadores);  Conversores CC-CC;  Fontes de Alimentação CC;  Inversores.}

\def\PUDconteudo{ & INTRODUÇÃO: Conceitos de Eletrônica de potência e Industrial. 
\\ 
 & INTERRUPTORES ELETRÔNICOS: Diodo.  Tiristores.  Transistor.  Escolha do interruptor.  Drivers para acionamento dos interruptores. 
\\ 
 & CÁLCULOS DE POTÊNCIA: Potência e Energia.  Potência Instantânea.  Energia.  Potência média.  Indutores e Capacitores.  Recuperação de Energia.  Valores Eficazes: RMS, Potência aparente e Fator de potência.  Cálculos de Potência para circuitos CA senoidais e não senoidais.  Cálculos de potência empregando o Simulador. 
\\ 
 & RETIFICADORES DE MEIA-ONDA: Carga Resistiva.  Carga Resistiva Indutiva.  Simulação de retificadores.  Fonte com carga RL.  Operação com diodo de roda livre.  Retificador de meia onda com filtro capacitivo.  Retificador de meia-onda Controlado.  Simulação de retificadores. 
\\ 
 & RETIFICADORES DE ONDA-COMPLETA: Operação com ponte e com tap-central.  Carga resistiva, carga RL.  Retificador de onda completa controlado.  Simulação de retificadores. 
\\ 
 & RETIFICADORES TRIFÁSICOS: Operação com carga resistiva.  carga RL.  Retificadores trifásicos totalmente controlados e operação de retificadores de doze pulsos. 
\\ 
 & CONTROLADORES DE TENSÃO CA (Gradadores): Monofásicos com carga R e Carga RL.  Trifásicos com carga resistiva e Indutiva.  Controle de velocidade do motor de indução com Gradadores. 
\\ 
 & CONVERSORES CC-CC: Conversor Buck (abaixador) operando em modo de condução contínua.  Relações de tensões e correntes.  Ondulação da tensão de saída.  Conversor Boost (elevador) operando em modo de condução contínua.  relações de tensões e correntes.  Ondulação da tensão de saída.  
Conversor Buck-Boost operando em modo de condução contínua.  relações de tensões e correntes.  Ondulação da tensão de saída.  Conversores Intercalados.  Operação dos conversores no modo de condução descontínua: Buck e Boost.  Simulações dos conversores CC-CC. 
\\ 
 & CONVERSORES CC-CC: Fontes de Alimentação CC.  Conversor Flyback operando em MCC e em MCD.  Conversor Forward (direto).  Conversor Push-Pull.  Conversor CC-CC em meia-ponte e em ponte-completa.  Simulações de conversores CC-CC. 
\\ 
 & INVERSORES: Inversor com senoide modificada.  Controle de Amplitude e Harmônica.  Inversor em meia-ponte.  Inversor em ponte completa.  Operação com modulação PWM senoidal uni e bipolar.  Inversores Trifásicos.  Simulações de inversores. }

\def\PUDmetodo{Aulas expositórias que podem ser teóricas e/ou práticas, onde as práticas no laboratório serão marcadas no decorrer da disciplina.}

\def\PUDavalia{A avaliação será feita com aplicação de provas teóricas e/ou práticas, além da possibilidade de inclusão de trabalhos e seminários no decorrer da disciplina.}

\def\PUDbibbasica{ & \href{www.biblioteca.ifce.edu.br/index.asp?codigo_sophia=23855}{Hart}, Daniel W.  Eletrônica de potência: análise e projetos de circuitos.  Porto Alegre, RS: AMGN, 2012.  479p.  ISBN: 8587918036.
\\ 
 & \href{www.biblioteca.ifce.edu.br/index.asp?codigo_sophia=24365}{Arrabaça}, Devair Aparecido; Gimenez, Salvador Pinilos.  Eletrônica de Potência Conversores de energia CA/CC.  São Paulo, SP: Érica, 2011.  334p.  ISBN: 9788536503714.
\\ 
 & \href{www.biblioteca.ifce.edu.br/index.asp?codigo_sophia=24217}{Mello}, Luiz Fernando Pereira.  Projetos de Fontes Chaveadas - Teoria e Prática.  São Paulo, SP: Érica, 2011.  284p.  ISBN: 9788536503370. }

\def\PUDbibcomplementar{ & \href{http://www.biblioteca.ifce.edu.br/index.asp?codigo_sophia=61869}{Rashid}, Muhhamad H.  Eletrônica de potência.  E-book. 4º ed.  Pearson.  884p.  ISBN: 9788543005942.
\\
 & \href{www.biblioteca.ifce.edu.br/index.asp?codigo_sophia=24053}{Franchi}, C. M.  Inversores de Frequência: Teoria e Aplicações.  2ª ed.  São Paulo, SP: Érica, 2011.  192p.  ISBN: 9788536502106.
\\ 
 & \href{www.biblioteca.ifce.edu.br/index.asp?codigo_sophia=24551}{Almeida}, José Luiz Antunes.  Dispositivos Semicondutores: Tiristores - Controle de potência em CC e CA.  12ª ed.  São Paulo, SP: Érica, 2013.  150p.  ISBN: 9788571942981.
\\
 & \href{www.biblioteca.ifce.edu.br/index.asp?codigo_sophia=23917}{Ahmed}, A.  Eletrônica de Potência.  São Paulo, SP: Pearson Prentice Hall, 2006.  479p.  ISBN: 8587918036.
\\
 & \href{www.biblioteca.ifce.edu.br/index.asp?codigo_sophia=24025}{Rodrigues}, Marcelo.  Gestão da manutenção elétrica, eletrônica e mecânica.  Curitiba, PR: Base Editorial, 2010.  128p.  ISBN: 9788579055690. }
%\\ 
 %&  RASHID, Muhammad H.  Eletrônica de Potência – circuitos, dispositivos e aplicações.  Makron Books, 1999.  ISBN 9788534605984
%\\ 
 %&  MARTINS, Denizar Cruz e BARBI, Ivo.  Conversores CC-CC Básicos não isolados.  UFSC, 2000.  ISBN 9788590104636. }

\def\PUDautor{Luiz Daniel Santos Bezerra}

\def\PUDdata{2013-04-23}

\def\PUDrevisao{1}

\def\PUDrevisor{Luiz Daniel Santos Bezerra}

\def\PUDdatarevisao{2019-05-15}

\PUDcorpo
